\section{Automated Screening}
\label{sec:automated-screening}

Organisations often rely on automated screening systems to filter, rank or manage large volumes of candidates.
These screening systems may also be referred to as Applicant Tracking Systems (ATS) throughout the rest of this report
for brevity. According to Jobscan's 2022 entry, we have that 97\% of large firms leverage these ATS platforms.
While there are no statistics on the adoption of small to medium-sized firms, it is reasonable to assume that
most organisations use ATS platforms to some extent due to the scalability and efficiency benefits they provide.

ATS platforms can be categorised in various ways, but in our report we will categorise them based on the underlying
logic of the scoring system. More specifically, we have Rule-Based, Keyword-Matching, and AI-Assisted ATS.

\subsection{Rule-Based ATS}
Rule-Based ATS systems refer to platforms that focus on explicit, pre-defined rules or criteria in order to 
screen candidates. Recruiters often configure such systems with keyword lists, required thresholds for
skills, education and experience, and other boolean queries. As such, these systems are knockout in
nature. Either a candidate satisfies every rule, or they are filtered out. It is important to note that
such systems do not provide a match score, nor do they rank candidates, instead recruiters are presented with
pass/fail shortlists. Due to their consistency and predictability, employers may leverage rule-based systems to screen candidates. 
By applying explicit, pre-defined criteria uniformly, such systems avoid the intuitive, subjective judgments that 
often influence human reviewers \cite{highhouse2008stubborn}.

The reason these systems are not commonly used however, is due to the limit of standardisation of CV formats.
These systems will struggle with non-standard CV formats, leading to parsing errors \cite{emnlp2025resumebench}.
Candidates with unconventional layouts or career paths are therefore more likely to be scored poorly for formatting 
reasons rather than skill.
Along with this, these systems are vulnerable to keyword-gaming, a practice where applicants strategically insert 
repeated or irrelevant terms to artificially inflate their match scores, a practice that is widely discussed in practitioner advice on optimising CVs for ATS parsers \cite{marr2019}.
The failure to recognise all possible CV formats and lack of semantic understanding also leads to the system being unable
to recognise unconventional experience. It may also perform poorly when skills are also presented in non-traditional ways, or
demonstrated in another section, such as open-source contributions or bootcamp projects. Rule-based systems are likely not 
exhaustive enough in the treatment of CVs and will likely score non-traditional CVs poorly, not due to the candidate's skills,
but instead due to the system's inability to recognise the candidate's skills.


\subsection{Keyword-Matching ATS}
This version of ATS is a middle ground between the previously mentioned rule-based approach, and the upcoming AI-driven version.
This form relies heavily on exact keyword matching and basic tokenisation. While some commercial versions may incorporate basic 
NLP-based parsing to recognise simple linguistic variations, such as standardising acronyms, they still offer limited
semantic understanding. They act more as advanced string-matching systems rather than engines that can understand the 
context of a candidate's experience.
Unlike rule-based knockouts, these systems typically assign match scores and rank candidates across the applicant pool rather than returning pass/fail shortlists alone.

Their strengths are mainly that they are simpler to implement and offer an increase in the transparency of the system compared 
to AI-driven systems, since much of the logic is still tied to explicit keyword counts.

Due to the simplicity of the system, it is often subject to misinterpretation issues, often failing to recognise
contextual details and ambiguous language within a CV. They also may overlook crucial nuances such as depth of skill,
quality of experience and also evidence of measurable achievement.
The misinterpretation then leads to misclassification. This is backed by employer surveys which report that more than 
90\% of employers use applicant-tracking or recruitment-management systems to filter or rank candidates, and 88--94\% 
of them believe that qualified applicants are screened out when they fail to match job-description criteria 
exactly \cite{fuller2021hiddenworkers}.

\subsection{AI-Assisted ATS}
AI-Assisted ATS comprise machine learning models, often leveraging transformers and Large Language Models (LLMs), 
to predict candidate-job fit. The system does not follow that of its predecessors, instead inferring semantic 
relationships between skills, roles and qualifications. Furthermore, many commercial versions will train on previous 
hiring data to learn which traits led to successful hires in the past.
Like keyword-matching ATS, they score and rank candidates across the applicant pool rather than applying knockout filters alone.

Due to this, the systems employ vector-embeddings and semantic search to evaluate whether a candidate is a strong 
match for the job description, or whether they are similar to other successful candidates for that role. They also take advantage
of predictive scoring models that learn from recruiter usage patterns to refine the scoring system,
picking up on hidden non-obvious signals that may be missed by traditional keyword-matching systems.

Given that the underlying system is well-trained, the key strength of AI-Assisted ATS lies in their
matchmaking accuracy, as well as detection of transferable skills. They are also able to handle 
a wider range of career paths compared to rule-based systems.

As mentioned before, the systems train on both historical hiring data and recruiter usage patterns.
This leads to a major flaw in these systems, the potential to reinforce historical biases.
Algorithmic bias is a concept where models inherit preference towards certain patterns, in this case,
it would be gender, ethnic or institutional preference found in the past hiring data \cite{raghavan2020,bogen2018}.
The potential for further discrimination is also present due to the feedback loop from recruiter interactions
if these employers continue to favour candidates with these specific traits.

Furthermore, the “black box” nature of these systems leads to low transparency, as there 
is often little explanation for ranking and rejecting decisions, undermining both 
candidate trust as well as compliance with fair-hiring guidelines.
Another significant issue lies in over-generalisation, where these models are often 
trained on broad datasets, prioritising generic markers over role-specific expertise.
In a similar fashion, domain drift also plays a role, where the system continues to 
optimise for outdated historical traits despite a company's technological or strategic shifts.

\vspace{1em}
Throughout this section, we have considered the different strengths and weaknesses of these ATS platforms, however,
there is still one major flaw that has not been discussed. This is the exclusive reliance on the candidate CV, 
completely disregarding all other sources of evidence that are typically provided by candidates (primarily GitHub and LinkedIn).
These systems work on a small subset of the data that is actually available to them, limiting their ability to truly capture
the strengths and weaknesses of a given candidate.

The sections that follow review research and technical foundations along three themes, that being configurable screening, 
multi-source verification, and explainable scoring. We then summarise how MERIT~\cite{merit_repo} is positioned 
to address these limitations (Section~\ref{sec:relevance-to-merit}).